% !TEX root = ../Main.tex
\chapter{练习题}

\section{填空题}

\ex{
    已知点 $O$ 为圆锥 $PO$ 底面的圆心，圆锥 $PO$ 的轴截面为边长为 $2$ 的等边三角形 $PAB$，求圆锥 $PO$ 的外接球的表面积。
}

\sol{
    轴截面是边长 $2$ 的等边三角形，所以底面半径 $r=1$、高 $h=\sqrt3$。球心在轴 $PO$ 上，套垂线模型 $R=\dfrac{h^2+r^2}{2h}$：
    \[
    R=\frac{3+1}{2\sqrt3}=\frac{2}{\sqrt3},\qquad R^2=\frac43.
    \]
    表面积 $=4\pi R^2=\dfrac{16\pi}{3}$。
}

\ex{
    已知三棱锥 $S$-$ABC$ 的所有顶点都在球 $O$ 的球面上，$SA\perp$ 平面 $ABC$，$SA=2\sqrt3$，$AB=1$，$AC=2$，$\angle BAC=\dfrac{\pi}{3}$，求球 $O$ 的表面积。
}

\sol{
    $SA$ 垂直底面，是侧棱垂直模型。先求底面外接圆半径：余弦定理
    \[
    BC^2=1+4-2\cdot1\cdot2\cdot\cos\tfrac{\pi}{3}=3,\quad BC=\sqrt3;
    \]
    正弦定理 $\dfrac{BC}{\sin\angle BAC}=2r$ 得 $r=\dfrac{\sqrt3}{2\cdot\frac{\sqrt3}{2}}=1$。套公式
    \[
    R^2=r^2+\Big(\frac{SA}{2}\Big)^2=1+3=4.
    \]
    表面积 $=4\pi R^2=16\pi$。
}

\ex{
    已知点 $A$ 是以 $BC$ 为直径的圆 $O$ 上异于 $B,C$ 的动点，$P$ 为平面 $ABC$ 外一点，且平面 $PBC\perp$ 平面 $ABC$，$BC=3$，$PB=2\sqrt2$，$PC=\sqrt5$，求三棱锥 $P$-$ABC$ 外接球的表面积。
}

\sol{
    两侧面 $PBC$、$ABC$ 共边 $BC$ 且互相垂直，$l=BC=3$。

    $\triangle ABC$ 的外接圆：$A$ 在以 $BC$ 为直径的圆上，所以这个圆就是它的外接圆，$r_1=\dfrac32$。

    $\triangle PBC$ 的外接圆：
    \[
    \cos\angle BPC=\frac{PB^2+PC^2-BC^2}{2\cdot PB\cdot PC}=\frac{8+5-9}{2\cdot2\sqrt2\cdot\sqrt5}=\frac{1}{\sqrt{10}},
    \quad \sin\angle BPC=\frac{3}{\sqrt{10}};
    \]
    \[
    r_2=\frac{BC}{2\sin\angle BPC}=\frac{3}{2\cdot\frac{3}{\sqrt{10}}}=\frac{\sqrt{10}}{2}.
    \]
    套公式
    \[
    R^2=r_1^2+r_2^2-\Big(\frac l2\Big)^2=\frac94+\frac{10}{4}-\frac94=\frac{5}{2}.
    \]
    表面积 $=4\pi R^2=10\pi$。
}

\ex{
    棱长为 $1$、各面都为等边三角形的四面体内有一点 $P$，由点 $P$ 向各面作垂线，垂线段的长度分别为 $d_1,d_2,d_3,d_4$，求 $d_1+d_2+d_3+d_4$。
}

\sol{
    正四面体各面面积相等，记为 $S$。以 $P$ 为顶点、各面为底分割：
    \[
    V=\frac13 S d_1+\frac13 S d_2+\frac13 S d_3+\frac13 S d_4=\frac{S}{3}(d_1+d_2+d_3+d_4).
    \]
    又以一个面为底、对顶点为顶：$V=\dfrac13 S h$。两式相等得
    \[
    d_1+d_2+d_3+d_4=h=\frac{\sqrt6}{3}\cdot1=\frac{\sqrt6}{3}.
    \]
    （内部任一点到各面距离之和都等于这个高，与点 $P$ 的位置无关。）
}

\ex{
    正三棱柱 $ABC$-$A_1B_1C_1$ 中，$AB=4$，$AA_1=6$，$E,F$ 分别是棱 $BB_1,CC_1$ 上的点，求三棱锥 $A$-$A_1EF$ 的体积。
}

\sol{
    $E,F$ 在与 $AA_1$ 平行的两条侧棱上，沿棱滑动不改变体积。取 $A=(0,0,0)$，$B=(4,0,0)$，$C=(2,2\sqrt3,0)$，$A_1=(0,0,6)$，$E=(4,0,t)$，$F=(2,2\sqrt3,s)$：
    \[
    V=\frac16\big|\det[\vec{AA_1},\vec{AE},\vec{AF}]\big|
    =\frac16\left|\det\begin{pmatrix}0&0&6\\4&0&t\\2&2\sqrt3&s\end{pmatrix}\right|.
    \]
    按第一行展开：$\det=6\cdot\det\begin{pmatrix}4&0\\2&2\sqrt3\end{pmatrix}=6\cdot8\sqrt3=48\sqrt3$，与 $t,s$ 无关。
    \[
    V=\frac{48\sqrt3}{6}=8\sqrt3.
    \]
}

\ex{
    底面边长为 $3,4,5$、高为 $6$ 的直三棱柱形容器内放置一气球，使气球尽可能膨胀（保持球的形状），求气球表面积的最大值。
}

\sol{
    $3,4,5$ 是直角三角形，其内切圆半径 $r_{\text{底}}=\dfrac{3+4-5}{2}=1$。球受两个限制：
    \[
    \text{径向：}\ R\le r_{\text{底}}=1;\qquad \text{轴向：}\ 2R\le6\ \Rightarrow\ R\le3.
    \]
    取更严的 $R=1$，表面积最大值 $=4\pi R^2=4\pi$。
}

\section{选择题}

\ex{
    如图，圆形纸片的圆心为 $O$，半径为 $6\,$cm，该纸片上的正方形 $ABCD$ 的中心为 $O$，$E,F,G,H$ 为圆 $O$ 上的点，$\triangle ABE,\triangle BCF,\triangle CDG,\triangle ADH$ 分别是以 $AB,BC,CD,DA$ 为底边的等腰三角形。沿虚线剪开后，分别以 $AB,BC,CD,DA$ 为折痕折起 $\triangle ABE,\triangle BCF,\triangle CDG,\triangle ADH$，使 $E,F,G,H$ 重合得到一个四棱锥。当该四棱锥的侧面积是底面积的 $2$ 倍时，该四棱锥的外接球的表面积为（\quad）。
    \[(A)\ \frac{16\pi}{3}\qquad (B)\ \frac{25\pi}{3}\qquad (C)\ \frac{64\pi}{3}\qquad (D)\ \frac{100\pi}{3}\]
}

\sol{
    设正方形边长 $a$。正方形中心 $O$ 到边 $AB$ 的距离是 $\dfrac a2$，顶点 $E$ 沿这条方向落在圆上，所以等腰三角形 $ABE$ 的高（斜高）$=6-\dfrac a2$。
    \[
    \text{底面积}=a^2,\qquad \text{侧面积}=4\cdot\frac12\,a\Big(6-\frac a2\Big)=12a-a^2.
    \]
    侧 $=2\times$ 底：$12a-a^2=2a^2\Rightarrow a=4$。

    折起后：底面正方形边长 $4$，半对角线 $2\sqrt2$；斜高 $=6-2=4$，底面中心到边的距离 $=2$，故顶点高
    \[
    h=\sqrt{4^2-2^2}=2\sqrt3.
    \]
    球心在轴上、距底面 $k$，由 $R^2=(2\sqrt3-k)^2=k^2+(2\sqrt2)^2$：
    \[
    12-4\sqrt3\,k=8\Rightarrow k=\frac{1}{\sqrt3},\qquad R^2=k^2+8=\frac13+8=\frac{25}{3}.
    \]
    表面积 $=4\pi R^2=\dfrac{100\pi}{3}$，选 (D)。
}

\ex{
    正方体的外接球与内切球上各有一个动点 $M,N$，若线段 $MN$ 的最小值为 $\sqrt3-1$，则（\quad）。
    \[(A)\ \text{正方体的外接球的表面积为 }12\pi\qquad (B)\ \text{正方体的内切球的体积为 }\frac{\pi}{3}\]
    \[(C)\ \text{正方体的棱长为 }1\qquad\qquad (D)\ \text{线段 }MN\text{ 的最大值为 }\sqrt3+1\]
}

\sol{
    设棱长 $a$。两球同心：内切球半径 $\dfrac a2$，外接球半径 $\dfrac{\sqrt3}{2}a$。$M,N$ 到公共球心的距离分别为 $\dfrac{\sqrt3}{2}a$、$\dfrac a2$，故
    \[
    \frac{\sqrt3-1}{2}a\le|MN|\le\frac{\sqrt3+1}{2}a.
    \]
    由 $\dfrac{\sqrt3-1}{2}a=\sqrt3-1$ 得 $a=2$。逐项检验：
    \begin{itemize}
        \item (A) 外接球半径 $\sqrt3$，表面积 $4\pi\cdot3=12\pi$。\checkmark
        \item (B) 内切球半径 $1$，体积 $\dfrac{4\pi}{3}\ne\dfrac{\pi}{3}$。$\times$
        \item (C) 棱长为 $2$。$\times$
        \item (D) $|MN|$ 最大值 $=\dfrac{\sqrt3+1}{2}\cdot2=\sqrt3+1$。\checkmark
    \end{itemize}
    选 (A)(D)。
}

\ex{
    已知三棱锥 $S$-$ABC$ 的所有顶点都在球 $O$ 的球面上，$\triangle ABC$ 是边长为 $1$ 的正三角形，$SC$ 为球 $O$ 的直径，$SC=2$，则此棱锥的体积为（\quad）。
    \[(A)\ \frac{\sqrt2}{6}\qquad (B)\ \frac{\sqrt3}{6}\qquad (C)\ \frac{\sqrt2}{3}\qquad (D)\ \frac{\sqrt2}{2}\]
}

\sol{
    $SC$ 是直径，$R=1$。球面上任一点对直径张直角，故 $\angle SAC=\angle SBC=90^\circ$，于是 $SA=\sqrt{SC^2-AC^2}=\sqrt{4-1}=\sqrt3$。

    求 $S$ 到平面 $ABC$ 的高。$\triangle ABC$ 外心 $G$ 满足 $GA=\dfrac{1}{\sqrt3}$；球心 $O$（即 $SC$ 中点）到平面 $ABC$ 的距离 $OG=\sqrt{R^2-GA^2}=\sqrt{1-\tfrac13}=\sqrt{\tfrac23}$。由 $O$ 是 $SC$ 中点、$C$ 在底面上，$S$ 到底面的距离 $=2\,OG=\dfrac{2\sqrt6}{3}$。
    \[
    V=\frac13\cdot\frac{\sqrt3}{4}\cdot\frac{2\sqrt6}{3}=\frac{\sqrt{18}}{18}=\frac{\sqrt2}{6}.
    \]
    选 (A)。
}

\ex{
    一个斜三棱柱的底面是边长为 $4$ 的正三角形，侧棱长 $5$，若其中一条侧棱与底面三角形的相邻两边都成 $60^\circ$ 角，则这个三棱柱的体积是（\quad）。
    \[(A)\ \frac{50\sqrt3}{3}\qquad (B)\ 20\sqrt3\qquad (C)\ \frac{50\sqrt2}{3}\qquad (D)\ 20\sqrt2\]
}

\sol{
    底面积 $=\dfrac{\sqrt3}{4}\cdot16=4\sqrt3$。

    设底面正三角形在顶点 $A$ 处的两边方向单位向量 $\vec u_1,\vec u_2$（$\vec u_1\cdot\vec u_2=\cos60^\circ=\tfrac12$），侧棱方向单位向量 $\vec v$ 与两者都成 $60^\circ$：$\vec v\cdot\vec u_1=\vec v\cdot\vec u_2=\tfrac12$。由对称性 $\vec v=p\,\dfrac{\vec u_1+\vec u_2}{|\vec u_1+\vec u_2|}+q\,\vec n$（$\vec n$ 为底面单位法向）。$|\vec u_1+\vec u_2|=\sqrt3$，由
    \[
    \vec v\cdot\vec u_1=p\cdot\frac{1+\frac12}{\sqrt3}=\frac{p\sqrt3}{2}=\frac12\ \Rightarrow\ p=\frac{1}{\sqrt3},\quad q^2=1-p^2=\frac23,\ q=\frac{\sqrt6}{3}.
    \]
    棱柱的高 $=$ 侧棱在法向的分量 $=5q=\dfrac{5\sqrt6}{3}$。
    \[
    V=4\sqrt3\cdot\frac{5\sqrt6}{3}=\frac{20\sqrt{18}}{3}=20\sqrt2.
    \]
    选 (D)。
}
